\subsection{Translations}

Let K be a (commutative) field, f a place of K, v a valuation on K and A a valuation ring of K. We shall say that A, f and v are associated if A is the ring of f and the ring of v. By virtue of no. 1 and § 2, no. 3, each of the three objects A, f and v then determines the other two (up to an equivalence as far as places and valuations are concerned). In particular there are the following equivalences:

\[
\begin{array}{l l l l} x \in \mathrm{A} & \Leftrightarrow f (x) \neq \infty & \Leftrightarrow v (x) \geqslant 0 \\ x \in \mathfrak {m} (\mathrm{A}) & \Leftrightarrow f (x) = 0 & \Leftrightarrow v (x) > 0 \\ x \in \mathrm{A} - \mathfrak {m} (\mathrm{A}) = \mathrm{U} (\mathrm{A}) \Leftrightarrow f (x) \neq 0 \text {and} f (x) \neq \infty \Leftrightarrow v (x) = 0 \\ x \in \mathrm{K} - \mathrm{A} & \Leftrightarrow f (x) = \infty & \Leftrightarrow v (x) <   0 \end{array}
\]

Every result relating to valuation rings, places or valuations can be translated into a result relating to the other two notions. Thus Proposition 4 of § 2, no. 4 gives:

PROPOSITION 5. Let K be a field, v a valuation on K and K' an extension of K. There exists a valuation v' on K' whose restriction to K is equivalent to v.

Let \(\Gamma_v\) and \(\Gamma_{v'}\) be the order groups of \(v\) and \(v'\) . Since the restriction of \(v'\) to K is equivalent to \(v\) , there exists an isomorphism \(\lambda\) of \(\Gamma_v\) onto a subgroup of \(\Gamma_{v'}\) , such that \(v' = \lambda \circ v\) on K. If \(\Gamma_v\) is identified with \(\lambda(\Gamma_v)\) by means of \(\lambda\) , it is seen that \(v'\) extends \(v\) .

Note that \(\Gamma_{v'}\) is in general distinct from \(\lambda(\Gamma_v)\) and the equivalence class of \(v'\) is not necessarily unique. We shall return to this in § 8.

Translating Theorem 3 of § 1, no. 3 (or Proposition 6 of § 2, no. 5), we obtain:

PROPOSITION 6. Let K be a field, A a subring of K and x an element of K. For x to be integral over A, it is necessary and sufficient that every valuation on K which is positive on A be positive at x.

From now on, we shall in general leave to the reader the trouble of performing translations analogous to the above.

